% =====================================================================
%  corps/00-prise-en-main.tex
%  Mode d'emploi général du modèle (structure, métadonnées, options).
%  Ce chapitre est destiné à l'auteur du mémoire/thèse : supprimez-le une
%  fois votre document en place.
% =====================================================================
\chapter{Prise en main du modèle}
\label{chap:prise-en-main}

\begin{quote}\small
Ce chapitre explique comment utiliser le modèle. \textbf{Supprimez-le}
(ainsi que le chapitre~\ref{chap:exemples}) lorsque vous rédigez votre
propre document : retirez les lignes \code{\input{...}} correspondantes
dans \code{these.tex}.
\end{quote}

% ---------------------------------------------------------------------
\section{Compilation}

Le document se compile avec \textbf{XeLaTeX} (pour l'Unicode et les polices
OpenType), suivi de \code{biber} pour la bibliographie :

\begin{latexcode}
xelatex these
biber   these
xelatex these
xelatex these
\end{latexcode}

\noindent Sur Overleaf, réglez le compilateur sur \textbf{XeLaTeX}
(\emph{Menu $\rightarrow$ Settings $\rightarrow$ Compiler}). En local, \code{latexmk} lit le
fichier \code{latexmkrc} fourni et enchaîne tout automatiquement :

\begin{latexcode}
latexmk these
\end{latexcode}

% ---------------------------------------------------------------------
\section{Structure des fichiers}

Le document maître est \code{these.tex}. Il ne contient presque pas de
texte : il charge la configuration puis inclut les chapitres.

\begin{latexcode}
% =====================================================================
%  config/preambule.tex
%  Réglages généraux : classe, géométrie, paquets de base.
%  Ce fichier ne charge ni les polices, ni les langues, ni le style
%  visuel FTI : chacun de ces aspects a son propre fichier de config.
% =====================================================================

% ---------------------------------------------------------------------
%  Vérification du moteur : ce modèle est prévu pour XeLaTeX.
%  (fontspec + polyglossia + rendu Unicode natif de l'arabe, du CJK…)
% ---------------------------------------------------------------------
\usepackage{iftex}
\RequireXeTeX % stoppe la compilation avec un message clair si ce n'est pas XeTeX

% ---------------------------------------------------------------------
%  Encodage d'entrée : avec XeLaTeX tout est UTF-8 nativement.
%  (Pas de \usepackage[utf8]{inputenc} ni de fontenc : inutiles/nuisibles.)
% ---------------------------------------------------------------------

% ---------------------------------------------------------------------
%  Géométrie de la page
%  typearea est déjà chargé par scrbook ; on ne le recharge pas (cela
%  provoquerait un « option clash »). On règle l'empagement via KOMAoptions.
%    BCOR = correction de reliure (marge intérieure supplémentaire).
%    DIV  = finesse de l'empagement KOMA (plus grand = marges + petites).
% ---------------------------------------------------------------------
\KOMAoptions{DIV=12,BCOR=8mm}
\recalctypearea

% ---------------------------------------------------------------------
%  Paquets de base
% ---------------------------------------------------------------------
\usepackage{xcolor}            % couleurs (utilisé par style-fti.tex)
\usepackage{graphicx}          % \includegraphics (logo, figures)
% microtype améliore la micro-typographie mais est facultatif : on ne
% le charge que s'il est installé (paquet Fedora : texlive-microtype).
\IfFileExists{microtype.sty}{\usepackage{microtype}}{}
\usepackage{booktabs}          % \toprule \midrule \bottomrule
\usepackage{tabularx}          % colonnes de largeur automatique
\usepackage{ragged2e}          % \RaggedRight, justification fine
\usepackage{setspace}          % \onehalfspacing pour le corps de la thèse
\usepackage{enumitem}          % listes personnalisables
\usepackage{csquotes}          % \enquote{} — guillemets adaptés à la langue

% ---------------------------------------------------------------------
%  Compatibilité TeX Live 2026 (arabe dans les tableaux)
%  Le noyau LaTeX 2026 a retiré la « fake math » autour des structures
%  tabulaires. Or le paquet « bidi » (chargé par polyglossia pour l'arabe,
%  voir config/langues.tex) appelle encore \UseMathForPositioningText,
%  désormais indéfini : dès qu'un tableau (tabular/tabularx) apparaît, on
%  obtient « Undefined control sequence » — y compris sur Overleaf (TL2026).
%  Ce shim rétablit la macro. Il est défini AVANT le chargement de bidi et
%  reste inoffensif sur les versions antérieures de TeX Live.
%  Réf. : https://www.overleaf.com/blog/tex-live-2026-is-now-available
%         (section « polyglossia, xepersian, and array »).
\makeatletter
\providecommand\UseMathForPositioningText{\m@th}
\makeatother

% Interligne : 1,5 est l'usage courant pour une thèse.
\onehalfspacing

% ---------------------------------------------------------------------
%  Coupures de ligne et débordements (« Overfull \hbox »)
%  ---------------------------------------------------------------------
%  Les noms d'auteurs en petites capitales issus de biblatex (p. ex.
%  « Schlechtweg ») et les liens hyperref ne se coupent pas toujours en
%  fin de ligne, ce qui peut faire déborder le texte dans la marge.
%  - \emergencystretch autorise TeX à relâcher l'espacement en dernier
%    recours plutôt que de déborder ;
%  - on autorise la césure à l'intérieur des liens hyperref (citations).
\emergencystretch=3em
\hfuzz=1pt
\hbadness=1500
% Autorise le retour à la ligne dans les hyperliens (citations, URL).
\usepackage{etoolbox}
\appto\UrlBreaks{\do\-}


% ---------------------------------------------------------------------
%  Titres et table des matières (réglages KOMA)
% ---------------------------------------------------------------------
\setcounter{tocdepth}{2}       % profondeur de la table des matières
\setcounter{secnumdepth}{3}    % profondeur de la numérotation des titres

% ---------------------------------------------------------------------
%  Liens hypertextes et métadonnées PDF.
%  hyperref doit être chargé tard ; polyglossia/bidi sont chargés
%  dans langues.tex, avant l'appel à ce bloc dans these.tex.
% ---------------------------------------------------------------------
\usepackage[
  unicode=true,                % métadonnées PDF en Unicode (titres non latins OK)
  hidelinks,                   % pas de cadres colorés autour des liens
  pdfusetitle,                 % titre/auteur PDF depuis \title/\author
  bookmarksnumbered=true,
  breaklinks=true,             % autorise un lien (citation, URL) à passer à la ligne
]{hyperref}

% Signets PDF corrects même avec des caractères Unicode.
\usepackage{bookmark}
   % classe, géométrie, paquets
% =====================================================================
%  config/polices.tex
%  Configuration typographique (fontspec — XeLaTeX).
%
%  Choix : polices 100 % LIBRES, déjà présentes dans TeX Live / la
%  distribution de polices Noto. Aucune police propriétaire.
%
%  Remarque sur la charte UNIGE : la police institutionnelle officielle
%  est « TheSans » (LucasFonts), qui est PROPRIÉTAIRE et ne peut pas être
%  redistribuée dans ce modèle. On utilise donc « Noto Sans » comme
%  substitut libre pour les éléments de type titre / logotype. Si vous
%  disposez d'une licence TheSans installée sur votre système, voir la
%  note « SUBSTITUTION THESANS » plus bas.
% =====================================================================

\usepackage{fontspec}
\defaultfontfeatures{Ligatures=TeX} % «--», «---», guillemets… gérés correctement

% ---------------------------------------------------------------------
%  1) POLICE PRINCIPALE (corps du texte)
%     Noto Serif : couverture Latin + Cyrillique + Grec (utile pour le
%     russe et le grec en écriture « inline » sans changer de police).
%
%  Repli automatique (voir la NOTE sur les polices variables plus bas et
%  scripts/get-fonts.sh) : si des faces STATIQUES de Noto Serif ont été
%  déposées dans ressources/polices/, on les charge par chemin ; sinon
%  on charge « Noto Serif » par son nom (cas normal sur Overleaf).
% ---------------------------------------------------------------------
\IfFileExists{ressources/polices/NotoSerif-Regular.ttf}{%
  \setmainfont{NotoSerif}[
    Path            = ressources/polices/,
    Extension       = .ttf,
    UprightFont     = *-Regular,
    BoldFont        = *-Bold,
    ItalicFont      = *-Italic,
    BoldItalicFont  = *-BoldItalic,
  ]%
}{%
  \setmainfont{Noto Serif}%
}

% ---------------------------------------------------------------------
%  2) POLICE SANS EMPATTEMENT (titres, page de titre, logotype)
%     Substitut libre de TheSans. Même repli statique que ci-dessus.
% ---------------------------------------------------------------------
\IfFileExists{ressources/polices/NotoSans-Regular.ttf}{%
  \setsansfont{NotoSans}[
    Path            = ressources/polices/,
    Extension       = .ttf,
    UprightFont     = *-Regular,
    BoldFont        = *-Bold,
    ItalicFont      = *-Italic,
    BoldItalicFont  = *-BoldItalic,
  ]%
}{%
  \setsansfont{Noto Sans}%
}

% ---------------------------------------------------------------------
%  3) POLICE À CHASSE FIXE (code, chemins, exemples techniques)
% ---------------------------------------------------------------------
\IfFileExists{ressources/polices/NotoSansMono-Regular.ttf}{%
  \setmonofont{NotoSansMono}[
    Path=ressources/polices/, Extension=.ttf,
    UprightFont=*-Regular, BoldFont=*-Bold,
    Scale=MatchLowercase,
  ]%
}{%
  \setmonofont{Noto Sans Mono}[Scale=MatchLowercase]%
}

% ---------------------------------------------------------------------
%  SUBSTITUTION THESANS (optionnel)
%  Si vous possédez TheSans, commentez le \setsansfont ci-dessus et
%  décommentez :
%
%  \setsansfont{TheSans}[
%    UprightFont = TheSans-Plain,
%    BoldFont    = TheSans-Bold,
%    ItalicFont  = TheSans-PlainItalic,
%  ]
% ---------------------------------------------------------------------

% =====================================================================
%  POLICES POUR LES ÉCRITURES NON LATINES
%  Déclarées comme familles séparées ; polyglossia (langues.tex) les
%  associe aux langues correspondantes. On garde aussi des commandes
%  \...font manuelles pour un usage direct.
% =====================================================================

% =====================================================================
%  NOTE IMPORTANTE SUR LES POLICES CJK ET ARABE ET LE PDF
%  ---------------------------------------------------------------------
%  XeLaTeX produit le PDF via « xdvipdfmx », qui NE SAIT PAS embarquer
%  une police VARIABLE (fichiers du type « NotoSerifCJK-VF.ttc » ou
%  « NotoNaskhArabic[wght].ttf »). Si votre système ne fournit QUE des
%  versions variables (cas de certaines distributions Linux récentes,
%  p. ex. Fedora), la compilation échoue avec :
%        xdvipdfmx:fatal: Invalid font
%
%  Sur OVERLEAF et la plupart des installations TeX Live, les versions
%  STATIQUES de Noto sont présentes : le chargement par NOM ci-dessous
%  fonctionne directement. Voir le README (section « Polices ») pour le
%  repli sur les systèmes à polices variables.
% =====================================================================

% --- Arabe (écriture de droite à gauche) -----------------------------
%  Noto Naskh Arabic : style naskh classique, très lisible.
%  Alternatives libres : « Amiri », « Scheherazade New ».
%
%  Repli automatique : si des instances STATIQUES ont été déposées dans
%  ressources/polices/ (voir scripts/get-fonts.sh et le README), on les
%  charge en priorité ; sinon on charge la police système par son nom.
\IfFileExists{ressources/polices/NotoNaskhArabic-Regular.ttf}{%
  \newfontfamily\arabicfont[
    Path        = ressources/polices/,
    Extension   = .ttf,
    UprightFont = NotoNaskhArabic-Regular,
    BoldFont    = NotoNaskhArabic-Bold,
    Script      = Arabic,
  ]{NotoNaskhArabic}%
}{%
  \newfontfamily\arabicfont{Noto Naskh Arabic}[Script=Arabic]%
}

% --- Chinois (simplifié) ---------------------------------------------
\IfFileExists{ressources/polices/NotoSerifCJKSC-Regular.otf}{%
  \newfontfamily\chinesefont[
    Path        = ressources/polices/,
    Extension   = .otf,
    UprightFont = NotoSerifCJKSC-Regular,
    BoldFont    = NotoSerifCJKSC-Bold,
  ]{NotoSerifCJKSC}%
}{%
  \newfontfamily\chinesefont{Noto Serif CJK SC}%
}

% --- Japonais --------------------------------------------------------
\IfFileExists{ressources/polices/NotoSerifCJKJP-Regular.otf}{%
  \newfontfamily\japanesefont[
    Path        = ressources/polices/,
    Extension   = .otf,
    UprightFont = NotoSerifCJKJP-Regular,
    BoldFont    = NotoSerifCJKJP-Bold,
  ]{NotoSerifCJKJP}%
}{%
  \newfontfamily\japanesefont{Noto Serif CJK JP}%
}

% --- Grec et Russe / Cyrillique --------------------------------------
%  Noto Serif couvre déjà le grec et le cyrillique ; on définit des
%  familles explicites pour que polyglossia dispose de \greekfont et
%  \russianfont dédiés. Même repli statique que la police principale.
\IfFileExists{ressources/polices/NotoSerif-Regular.ttf}{%
  \newfontfamily\greekfont[
    Path=ressources/polices/, Extension=.ttf,
    UprightFont=*-Regular, BoldFont=*-Bold,
    ItalicFont=*-Italic, BoldItalicFont=*-BoldItalic,
  ]{NotoSerif}%
  \newfontfamily\russianfont[
    Path=ressources/polices/, Extension=.ttf,
    UprightFont=*-Regular, BoldFont=*-Bold,
    ItalicFont=*-Italic, BoldItalicFont=*-BoldItalic,
  ]{NotoSerif}%
}{%
  \newfontfamily\greekfont{Noto Serif}%
  \newfontfamily\russianfont{Noto Serif}%
}
     % fontspec (latin, CJK, arabe, éthiopien)
% =====================================================================
%  config/langues.tex
%  Gestion multilingue avec polyglossia (adapté à XeLaTeX).
%
%  Langue principale : FRANÇAIS (typographie française : espaces
%  insécables avant : ; ! ?, guillemets « … », etc.).
%
%  Langues secondaires pré-configurées, avec leur écriture :
%    - arabe    (arabe,   RTL — vrai rendu droite-à-gauche via bidi)
%    - russe    (cyrillique)
%    - chinois  (hans/CJK)
%    - japonais (CJK)
%    - grec     (grec moderne/ancien)
%    - anglais  (pour l'abstract)
%
%  IMPORTANT : polyglossia doit être chargé APRÈS fontspec (polices.tex)
%  et les familles \arabicfont, \russianfont, etc. doivent déjà exister.
%
%  ARABE + TABLEAUX : l'option « arabic » charge le paquet « bidi », qui
%  assure un vrai algorithme de paragraphe RTL (alignement à droite,
%  césure, ordre des lignes). Sous TeX Live 2026 (dont Overleaf), bidi
%  provoque une erreur « \UseMathForPositioningText » dans les tableaux ;
%  le correctif est appliqué dans config/preambule.tex (shim), chargé
%  AVANT bidi. Voir ce fichier pour les détails.
% =====================================================================

\usepackage{polyglossia}

% --- Langue principale -----------------------------------------------
\setmainlanguage{french}

% --- Langues secondaires ---------------------------------------------
%  polyglossia associe automatiquement \arabicfont, \russianfont,
%  \chinesefont, \japanesefont, \greekfont (définis dans polices.tex).
\setotherlanguage{english} % pour l'abstract en anglais
\setotherlanguage{arabic}
\setotherlanguage{russian}
\setotherlanguage[variant=simplified]{chinese}
\setotherlanguage{japanese}
\setotherlanguage{greek}

% ---------------------------------------------------------------------
%  Commandes de confort pour insérer un court fragment dans une écriture
%  donnée, sans avoir à ouvrir un environnement complet.
%
%  Exemples d'utilisation :
%     \ar{النص العربي}     % arabe (RTL géré par polyglossia/bidi)
%     \ru{Русский текст}   % russe
%     \zh{中文文本}         % chinois
%     \ja{日本語のテキスト}  % japonais
%     \gr{Ἑλληνικά}        % grec
% ---------------------------------------------------------------------
\newcommand{\ar}[1]{\textarabic{#1}}
\newcommand{\ru}[1]{\textrussian{#1}}
\newcommand{\zh}[1]{\textchinese{#1}}
\newcommand{\ja}[1]{\textjapanese{#1}}
\newcommand{\gr}[1]{\textgreek{#1}}

% ---------------------------------------------------------------------
%  Passages longs
%  ---------------------------------------------------------------------
%  Utilisez les environnements fournis par polyglossia, qui gèrent la
%  césure et la direction du texte :
%     \begin{russian} ... \end{russian}
%     \begin{chinese} ... \end{chinese}   etc.
%
%  Pour un paragraphe ARABE, utilisez l'environnement « Arabic » avec une
%  MAJUSCULE (pour éviter le conflit avec la commande \arabic de LaTeX) :
%     \begin{Arabic} ... \end{Arabic}
%  Il assure un véritable rendu droite-à-gauche du bloc (alignement à
%  droite, retours à la ligne et césure corrects).
% ---------------------------------------------------------------------
     % polyglossia (français + langues secondaires)
% =====================================================================
%  config/biblio.tex
%  Bibliographie avec biblatex + biber (compatible Unicode/multilingue).
%  Style : auteur-date, courant en traductologie et sciences humaines.
% =====================================================================

\usepackage[
  backend=biber,          % biber : gère l'UTF-8 et le tri multilingue
  style=authoryear,       % (auteur, année)
  sorting=nyt,            % tri : nom, année, titre
  sortlocale=fr_FR,       % tri respectant les règles françaises
  language=french,        % langue de la bibliographie (libellés « In », etc.)
  autolang=none,          % pas de bascule auto : on gère les écritures non
                          % latines à la main via \ar{}, \zh{} (voir .bib).
                          % Évite les avertissements « Language not supported ».
  maxbibnames=99,         % liste tous les auteurs dans la bibliographie
  maxcitenames=2,         % « et al. » au-delà de 2 auteurs dans le texte
  giveninits=true,        % initiales des prénoms
  uniquename=init,
  doi=true, isbn=false, url=true,
]{biblatex}

% Fichier de références.
\addbibresource{biblio/references.bib}

% ---------------------------------------------------------------------
%  Note : au chargement, biblatex peut émettre quelques avertissements
%  « Language 'arabic'/'chinese'/'japanese' not supported » car il n'existe
%  pas de fichier de localisation (.lbx) pour ces langues déclarées à
%  polyglossia. Ces avertissements sont SANS CONSÉQUENCE : la bibliographie
%  est en français et les rares champs en écriture non latine sont gérés à
%  la main via \ar{} / \zh{} dans le fichier .bib.
%
%  IMPORTANT : ne PAS utiliser \DeclareLanguageMapping{chinese}{french}
%  (etc.) pour masquer ces avertissements : combiné à \setotherlanguage
%  {chinese}, cela casse la localisation française de biblatex (libellés
%  « In »/« byeditor » en anglais et sans espace). On préfère donc laisser
%  ces avertissements bénins.

% ---------------------------------------------------------------------
%  Réglages fins conseillés.
% ---------------------------------------------------------------------
% Guillemets français autour des titres d'articles (via csquotes).
\DeclareFieldFormat[article,inbook,incollection,inproceedings]{title}%
  {\enquote{#1}}

% Espace insécable avant les deux-points, etc. : géré par polyglossia
% (langue française) — rien à faire ici.
      % biblatex + biber (APA 7)
% =====================================================================
%  config/style-fti.tex
%  Identité visuelle : couleurs de la charte graphique de l'Université
%  de Genève, en-têtes/pieds de page, mise en forme des titres.
%
%  Sources des couleurs (charte graphique UNIGE, Service de communication) :
%    https://www.unige.ch/communication/publier/charte/couleurs
%  Valeurs officielles :
%    - Rouge institutionnel UNIGE : Pantone 214C / #CF0063 / RVB 207,0,99
%    - Gris du sceau              : 45 % noir     / #A3A3A3 / RVB 163,163,163
%    - Couleur de la FACULTÉ de traduction et d'interprétation (FTI) :
%                                   Pantone 152C  / #FF5C00 / RVB 255,92,0
% =====================================================================

% ---------------------------------------------------------------------
%  Palette
% ---------------------------------------------------------------------
\definecolor{unigeRouge}{HTML}{CF0063} % rouge institutionnel (Pantone 214C)
\definecolor{unigeGris}{HTML}{A3A3A3}  % gris du sceau (45 % noir)
\definecolor{ftiOrange}{HTML}{FF5C00}  % couleur de la faculté FTI (Pantone 152C)

%  Couleur d'accent employée dans tout le document. Par défaut : la
%  couleur de la FTI (orange). Passez à \unigeRouge pour un rendu
%  « institutionnel » plus neutre.
\colorlet{accent}{ftiOrange}

% ---------------------------------------------------------------------
%  En-têtes et pieds de page (scrlayer-scrpage, KOMA)
% ---------------------------------------------------------------------
\usepackage[headsepline]{scrlayer-scrpage}
\pagestyle{scrheadings}
\clearpairofpagestyles

% Filet de séparation d'en-tête à la couleur d'accent.
\setkomafont{headsepline}{\color{accent}}
\KOMAoptions{headsepline=0.8pt}

% Contenu des en-têtes/pieds :
%   - en-tête extérieur : titre courant de chapitre/section
%   - pied extérieur    : numéro de page
\automark[section]{chapter}
\ihead{}
\ohead{\small\headmark}
\ofoot*{\normalfont\color{accent}\pagemark}
\cfoot*{}

% Marques de chapitre/section en style « propre » (sans « Chapitre N. »).
\renewcommand*{\chaptermarkformat}{}

% ---------------------------------------------------------------------
%  Mise en forme des titres (KOMA : \setkomafont / \addtokomafont)
% ---------------------------------------------------------------------
% Titres en police sans empattement, teintés à la couleur d'accent.
\addtokomafont{disposition}{\sffamily}
\addtokomafont{chapter}{\color{accent}}
\addtokomafont{section}{\color{accent}}
\addtokomafont{subsection}{\color{accent!85!black}}

% Numéro de chapitre imposant (style KOMA « chapterprefix » désactivé,
% numéro intégré au titre).
\renewcommand*{\chapterformat}{%
  \mbox{\textcolor{accent}{\Huge\thechapter}\enskip}%
}

% ---------------------------------------------------------------------
%  Liens hypertextes teintés à l'accent (discret, hyperref = hidelinks
%  masque les cadres ; ici on colore juste le texte des références).
% ---------------------------------------------------------------------
\hypersetup{
  colorlinks = false,   % pas de couleur agressive pour l'impression
  linkcolor  = accent,
  citecolor  = accent,
  urlcolor   = accent,
}

% ---------------------------------------------------------------------
%  Petites commandes utilitaires pour la page de titre / le style.
% ---------------------------------------------------------------------
% Filet horizontal à la couleur d'accent.
\newcommand{\ftirule}[1][\linewidth]{%
  {\color{accent}\rule{#1}{1.2pt}}%
}
   % couleurs de la charte, en-têtes, titres

% =====================================================================
%  corps/01-introduction.tex
% =====================================================================
\chapter{Introduction}
\label{chap:introduction}

Ce chapitre introduit la problématique de recherche, le contexte
scientifique et la structure de la thèse.

\section{Contexte et motivation}

La traductologie contemporaine s'appuie de plus en plus sur des
méthodes outillées et des corpus multilingues~\autocite{toury1995}.
Comme le rappelle \textcite{eco2003}, traduire revient souvent à
« dire presque la même chose ». Les travaux menés au sein du
Département TIM illustrent cette évolution, qu'il s'agisse de
traduction automatique de dialectes~\autocite{alalmaoui2025arabizi,mutal2026aladdin},
de traduction de la parole en milieu médical~\autocite{bouillon2021babeldr},
de terminologie et de variation~\autocite{humbertdroz2022higgs}, de
localisation~\autocite{morado2020xliff} ou encore de détection
automatique du changement sémantique~\autocite{tahmasebi2021semanticchange,schlechtweg2020semeval}.

\section{Question de recherche}

% Formulez ici votre question de recherche et vos hypothèses.
Nous posons la question suivante : \dots

\section{Structure de la thèse}

La thèse est organisée comme suit. Le
chapitre~\ref{chap:exemples} illustre les capacités multilingues du
présent modèle. La conclusion (chapitre~\ref{chap:conclusion})
synthétise les apports.

% =====================================================================
%  corps/02-exemples-multilingues.tex
%  Chapitre vitrine + MODE D'EMPLOI : chaque fonctionnalité est présentée
%  avec son code source (environnement « latexcode ») puis son rendu.
% =====================================================================
\chapter{Exemples multilingues et mode d'emploi}
\label{chap:exemples}

Ce chapitre illustre l'insertion de texte dans différentes écritures et
sert en même temps de \emph{mode d'emploi} : pour chaque fonctionnalité,
le code source est affiché, suivi de son rendu. Chaque langue secondaire
est pré-configurée dans \code{config/langues.tex} avec sa police adéquate
(\code{config/polices.tex}).

% ---------------------------------------------------------------------
\section{Fragments courts en ligne}

Le modèle définit des commandes de confort pour insérer un mot ou une
expression dans une écriture donnée. On les utilise ainsi (remplacez le
texte entre accolades par votre contenu dans l'écriture voulue) :

\begin{latexcode}
Arabe    : \ar{ ... texte arabe ... }      % droite à gauche
Russe    : \ru{ ... texte russe ... }
Chinois  : \zh{ ... texte chinois ... }
Japonais : \ja{ ... texte japonais ... }
Grec     : \gr{ ... texte grec ... }
Tigrinya : \ti{ ... texte tigrinya ... }
\end{latexcode}

\noindent \textbf{Rendu :}

\begin{itemize}
  \item Arabe : \ar{مرحبا بالعالم}
  \item Russe : \ru{Здравствуй, мир}
  \item Chinois : \zh{你好，世界}
  \item Japonais : \ja{こんにちは世界}
  \item Grec : \gr{Γειά σου, κόσμε}
  \item Tigrinya : \ti{ሰላም ዓለም}
\end{itemize}

Ces fragments cohabitent sans difficulté dans un paragraphe français, y
compris pour une langue de droite à gauche. Par exemple, le code

\begin{latexcode}
comme l'arabe \ar{ ... } qui reprend ensuite le fil du texte.
\end{latexcode}

\noindent produit : comme l'arabe \ar{النص العربي} qui reprend ensuite le
fil du texte français.

% ---------------------------------------------------------------------
\section{Passages longs (environnements)}

Pour un passage plus long, on utilise l'environnement de la langue, qui
gère la césure et la direction du texte. La syntaxe générale est :

\begin{latexcode}
\begin{Arabic}
  ... votre paragraphe en arabe ...
\end{Arabic}

\begin{russian}
  ... votre paragraphe en russe ...
\end{russian}
\end{latexcode}

\begin{quote}\small
\textbf{Attention.} L'environnement arabe s'écrit avec une \emph{majuscule}
(\code{Arabic}) pour éviter un conflit avec la commande \code{\arabic} de
LaTeX. Les autres s'écrivent en minuscules (\code{russian}, \code{chinese},
\code{japanese}, \code{greek}).
\end{quote}

\subsection{Arabe (droite à gauche)}

Un paragraphe complet en arabe est justifié et aligné à droite, avec un
ordre des lignes et une césure corrects (paquet \code{bidi}) :

\begin{Arabic}
تُسَاعِدُ اللُّغَاتُ المُخْتَلِفَةُ عَلَى التَّوَاصُلِ بَيْنَ النَّاسِ وَالثَّقَافَاتِ. وَمَعَ تَطَوُّرِ التِّقْنِيَاتِ الرَّقْمِيَّةِ، أَصْبَحَ مِنَ السَّهْلِ كِتَابَةُ النُّصُوصِ وَقِرَاءَتُهَا بِأَحْرُفٍ وَاتِّجَاهَاتٍ مُخْتَلِفَةٍ. وَيُسَاعِدُ دَعْمُ النُّصُوصِ مِنَ الْيَمِينِ إِلَى الْيَسَارِ عَلَى عَرْضِ اللُّغَةِ الْعَرَبِيَّةِ بِشَكْلٍ وَاضِحٍ وَطَبِيعِيٍّ.
\end{Arabic}

\subsection{Russe (cyrillique)}
\begin{russian}
Русский язык использует кириллический алфавит. Этот абзац демонстрирует
корректный набор кириллического текста, включая переносы и
типографику.
\end{russian}

\subsection{Chinois}
\begin{chinese}
这是一段中文示例文字，用于演示本模板对中日韩表意文字的排版支持。
中文通常不使用空格分词。
\end{chinese}

\subsection{Japonais}
\begin{japanese}
これは日本語のサンプル段落です。ひらがな、カタカナ、漢字が
混在しても正しく組版されることを示します。
\end{japanese}

\subsection{Grec}
\begin{greek}
Αὕτη ἡ παράγραφος εἶναι γραμμένη στὰ ἑλληνικά, γιὰ νὰ δείξει τὴν
ὑποστήριξη τοῦ ἑλληνικοῦ ἀλφαβήτου, μαζὶ μὲ τόνους καὶ πνεύματα.
\end{greek}

% ---------------------------------------------------------------------
\section{Tableaux (booktabs)}

Les tableaux se composent avec \code{booktabs} (filets \code{\toprule},
\code{\midrule}, \code{\bottomrule}) et \code{tabularx} (colonne \code{X}
de largeur automatique). Ils acceptent les commandes multilingues :

\begin{latexcode}
\begin{table}[htbp]
  \centering
  \begin{tabularx}{\linewidth}{@{}l X@{}}
    \toprule
    \textbf{Langue} & \textbf{« traduction »} \\
    \midrule
    Français & traduction \\
    Arabe    & \ar{ ... } \\
    Chinois  & \zh{ ... } \\
    \bottomrule
  \end{tabularx}
  \caption{Légende du tableau.}
  \label{tab:traduction}
\end{table}
\end{latexcode}

\noindent \textbf{Rendu :}

\begin{table}[htbp]
  \centering
  \begin{tabularx}{\linewidth}{@{}l X@{}}
    \toprule
    \textbf{Langue} & \textbf{« Traduction » (le mot pour « traduction »)} \\
    \midrule
    Français  & traduction \\
    Arabe     & \ar{ترجمة} \\
    Russe     & \ru{перевод} \\
    Chinois   & \zh{翻译} \\
    Japonais  & \ja{翻訳} \\
    Grec      & \gr{μετάφραση} \\
    Tigrinya  & \ti{ትርጉም}\\
    \bottomrule
  \end{tabularx}
  \caption{Le terme « traduction » dans plusieurs langues et écritures.}
  \label{tab:traduction}
\end{table}

% ---------------------------------------------------------------------
\section{Notes de bas de page et citations}

Une note de bas de page s'insère avec \code{\footnote{...}} et peut, elle
aussi, contenir plusieurs écritures :

\begin{latexcode}
\dots{} plusieurs écritures.\footnote{Arabe \ar{ ... }, russe \ru{ ... }.}
\end{latexcode}

\noindent \textbf{Rendu :} une note peut contenir plusieurs
écritures.\footnote{Exemple : arabe \ar{هامش}, russe \ru{сноска},
chinois \zh{脚注}, japonais \ja{脚注}.}

Les références bibliographiques se citent avec \code{\autocite{clé}}
(entre parenthèses) ou \code{\textcite{clé}} (dans la phrase) :

\begin{latexcode}
D'après \textcite{barhudarov1975}, la traduction \dots{}  \textcite{xietianzhen1999, aljahiz}
\dots{} la traductologie outillée \autocite{schlechtweg-etal-2020-semeval}.
\end{latexcode}

\noindent \textbf{Rendu :} D'après \textcite{barhudarov1975}, la traduction \dots{} \textcite{xietianzhen1999, aljahiz} 
\dots{} la traductologie outillée \autocite{schlechtweg-etal-2020-semeval}.

Comme démontré ci-dessus, les références à des auteurs dont le nom s'écrit dans des caractères non latins s'affichent et se
trient correctement grâce à \code{biber}.

\code{biber} permet également d'ajouter des \enquote{pre-notes} et \enquote{post-notes} :
\begin{latexcode}
\dots{} nous pouvons mentionner \textcite[originalement du 9ème siècle,][qui traduit Aristote]{aljahiz}
\end{latexcode}

\noindent \textbf{Rendu :} \dots{} nous pouvons mentionner \textcite[originalement du 9ème siècle,][qui traduit Aristote]{aljahiz}
% ... ajoutez vos chapitres ici ...
\end{latexcode}

Pour ajouter un chapitre, créez un fichier dans \code{corps/} et
ajoutez la ligne \code{\input{corps/mon-chapitre}} au bon endroit.

% ---------------------------------------------------------------------
\section{Renseigner les métadonnées}

Les informations de la page de titre se trouvent en haut de
\code{these.tex}. Modifiez-les :

\begin{latexcode}
\newcommand{\ThesisTitle}{Titre de la thèse}
\newcommand{\ThesisSubtitle}{Sous-titre éventuel}
\newcommand{\ThesisAuthor}{Prénom \textsc{Nom}}
\newcommand{\ThesisDegree}{docteur·e en traitement informatique multilingue}
\newcommand{\ThesisSupervisor}{Prof. Prénom \textsc{Nom}}
\newcommand{\ThesisJury}{Prof. A. \textsc{Un}, Université de Genève\\%
                         Prof. B. \textsc{Deux}, University of Stuttgart
                         }
\newcommand{\ThesisPlace}{Genève}
\newcommand{\ThesisDate}{2026}
\end{latexcode}

\noindent Les membres du jury sont séparés par \code{\\} (chacun sur sa
ligne). Ces champs alimentent aussi les métadonnées du PDF.

% ---------------------------------------------------------------------
\section{Couleur d'accent}

La couleur d'accent (titres, filets, liens) est définie dans
\code{config/style-fti.tex}. Par défaut, c'est l'orange de la FTI ; pour
un rendu institutionnel plus neutre, basculez sur le rouge UNIGE :

\begin{latexcode}
\colorlet{accent}{ftiOrange}   % défaut : couleur de la FTI (Pantone 152C)
% \colorlet{accent}{unigeRouge} % variante : rouge UNIGE (Pantone 214C)
\end{latexcode}

% ---------------------------------------------------------------------
\section{Structurer le texte}

Les commandes de structure usuelles sont disponibles :

\begin{latexcode}
\chapter{Titre de chapitre}
\section{Une section}
\subsection{Une sous-section}

\begin{itemize}
  \item premier point ;
  \item second point.
\end{itemize}

Une citation courte se met entre guillemets français avec
\enquote{le texte cité} (guillemets « … » automatiques).
\end{latexcode}

\noindent Pour insérer une figure :

\begin{latexcode}
\begin{figure}[htbp]
  \centering
  \includegraphics[width=0.8\linewidth]{ressources/mon-image}
  \caption{Légende de la figure.}
  \label{fig:mon-image}
\end{figure}
\end{latexcode}

On renvoie ensuite à un élément numéroté avec \code{\ref{}} (p. ex.
\code{\ref{fig:mon-image}}) et à une page avec \code{\pageref{}}.
